% =====================================================================
%  ADD-ON TEXT for "Building Personalized AI Judges"
%  All numbers seed-averaged; sweep = 30 raters x 14 held-out x 5 dims x 3 seeds
%  (run 20260731_134005_ctxsweep); full pool = 666 raters x 14 x 5 x 3 seeds
%  (report_ctx20/). car_mean baseline is leave-one-out.
% =====================================================================

% ---- (Replace / augment the existing methodology paragraph) ----
We evaluated OpenAI's GPT-5.6 Sol reasoning model (\texttt{gpt-5.6-sol}) through
the Responses API, using standard reasoning mode at medium reasoning effort; all
experiments were run in July--August 2026. To isolate the effect of context, we
swept the number of in-context exemplars over $N \in \{0, 4, 8, 12, 16, 20\}$
while holding every other lever fixed (a persona-style system prompt, exemplars
drawn uniformly at random from the rater's own remaining vehicles, and image-before-text
demonstration ordering). For each rater we reserved a fixed test set of $14$
vehicles from the $34$ they rated; the same held-out set was used at every context
size, so any two configurations are directly comparable cell-for-cell. Because the
deepest configuration ($N{=}20$) requires each rater to supply $20$ exemplars in
addition to the $14$ test vehicles, we admitted a vehicle's four shared anchor
ratings into the exemplar pool and retained only raters with enough densely rated
vehicles to satisfy the largest context size. We report mean absolute error (MAE),
signed bias (the mean of predicted minus human rating), single-rater absolute-agreement
ICC$(2,1)$, and Spearman's rank correlation $\rho$, all against each rater's own
held-out ratings. Owing to the cost of the reasoning model, the sweep used $30$
raters ($30 \times 14 \times 5 = 2{,}100$ evaluation cells per seed); pooling three
seeds gives $6{,}300$ paired cells clustered within $30$ raters. Having identified
$N{=}20$ as the strongest configuration, we re-ran it on the full eligible pool of
$666$ raters ($46{,}620$ cells per seed; ${\sim}140{,}000$ pooled over three seeds)
to obtain stable per-dimension estimates and to compare it against three
reference predictors: a \emph{population-mean} baseline (predict each dimension's
grand mean), a leave-one-out \emph{per-vehicle mean} baseline (predict the mean
rating other raters gave that vehicle), and the \emph{zero-shot} judge ($N{=}0$).

% =====================================================================
\subsubsection{Context Reduces Error, Corrects Calibration, and Improves Agreement}
\label{sec:res-contextsweep}
Adding context monotonically improved every metric, with the largest gains
concentrated in the first few exemplars (\Cref{fig:ch5_contextvsmae,fig:ch5_contextvsagreement}).
Overall MAE fell from $1.15$ at zero-shot to $1.01$ with only $N{=}4$ exemplars---a
$12\%$ reduction from a handful of examples---and continued to $0.93$ ($N{=}8$),
$0.88$ ($N{=}12$), $0.86$ ($N{=}16$), and $0.84$ ($N{=}20$). The marginal benefit
diminished past $N{=}12$ (each further step lowered MAE by only ${\sim}0.02$), but
$N{=}20$ remained the best configuration on \emph{every} metric, and we adopt it as
our judge. Agreement rose in tandem: ICC$(2,1)$ climbed from $0.44$ to $0.67$ and
Spearman $\rho$ from $0.46$ to $0.67$ over the same range, moving the judge from
weak into substantial agreement with individual raters.

The most striking effect of context is on \emph{calibration} rather than raw error
(\Cref{fig:ch5_contextvsbias}). At zero-shot the judge is systematically
miscalibrated, and the miscalibration is not uniform across dimensions. The model
\emph{over}-predicts every attribute and, especially, overall preference
(bias $+0.85$: it assumes people like the vehicles far more than they do), while
ruggedness is the \emph{only} dimension it \emph{under}-predicts (bias $-0.17$).
The remaining attributes sit in between (sporty $+0.18$, modern $+0.22$,
luxurious $+0.31$). Providing even a small amount of the rater's own context
collapses this bias almost entirely: by $N{=}4$ the overall absolute bias drops from
$0.28$ to $0.08$, the preference bias falls from $+0.85$ to $+0.23$, and
luxuriousness is nearly eliminated ($+0.31 \rightarrow +0.02$). In other words, the
exemplars teach the model \emph{where the person's scale sits}---their leniency and
per-attribute tendencies---which a generic zero-shot prompt cannot infer. Notably,
ruggedness remains the lone (slightly) negatively biased dimension at all context
sizes, foreshadowing its distinct behavior discussed below.

% =====================================================================
\subsubsection{The $N{=}20$ Judge Substantially Outperforms Non-Personalized Baselines}
\label{sec:res-ctx20}
On the full pool of $666$ raters, the $N{=}20$ judge attains an overall MAE of
$0.80$, ICC$(2,1)$ of $0.70$, and Spearman $\rho$ of $0.70$
(\Cref{fig:3seed_overall}), i.e. it reproduces an individual's held-out ratings
within roughly $0.8$ points on a $1$--$6$ scale and orders that person's vehicles
in substantial agreement with them. This markedly exceeds all three non-personalized
references (\Cref{tab:ctx20-improvement}). Relative to the population-mean baseline,
the judge reduces MAE by $37\%$ while lifting agreement from essentially no ranking
ability (ICC $0.13$, $\rho$ $0.24$) to substantial agreement (ICC $0.70$, $\rho$
$0.70$). It also improves on the far stronger per-vehicle-mean baseline---which
already encodes crowd consensus for each vehicle---by $27\%$ in MAE and by $56\%$
and $37\%$ in ICC and $\rho$ respectively, confirming that the personalized judge
captures individual taste that a consensus predictor cannot. Against the zero-shot
judge, personalization lowers MAE by $31\%$ and raises ICC by $63\%$. Every
comparison is significant under a rater-clustered bootstrap ($p < 0.001$).

\begin{table}[htbp]
\centering
\caption{Overall performance of the $N{=}20$ personalized judge and its
improvement over each non-personalized baseline (full pool, $666$ raters,
seed-averaged over three seeds). MAE improvement is a percent \emph{reduction};
ICC and $\rho$ improvements are percent \emph{increases}. The per-vehicle-mean
baseline is leave-one-out.}
\label{tab:ctx20-improvement}
\begin{tabular}{lccc}
\hline
 & MAE $\downarrow$ & ICC$(2,1)$ $\uparrow$ & Spearman $\rho$ $\uparrow$ \\
\hline
Population mean (constant)      & $1.267$ & $0.128$ & $0.243$ \\
Per-vehicle mean (leave-one-out)& $1.090$ & $0.447$ & $0.510$ \\
Zero-shot ($N{=}0$)             & $1.161$ & $0.427$ & $0.448$ \\
\textbf{Personalized judge ($N{=}20$)} & $\mathbf{0.797}$ & $\mathbf{0.698}$ & $\mathbf{0.697}$ \\
\hline
Improvement vs.\ population mean & $-37\%$ & $+446\%$ & $+187\%$ \\
Improvement vs.\ per-vehicle mean& $-27\%$ & $+56\%$  & $+37\%$  \\
Improvement vs.\ zero-shot       & $-31\%$ & $+63\%$  & $+55\%$  \\
\hline
\end{tabular}
\vspace{0.4em}
\begin{minipage}{\linewidth}
\footnotesize
\textit{Note}. Population-mean and per-vehicle-mean baselines are computed on the
identical held-out cells as the judge. The zero-shot values are from the matched
$N{=}0$ condition. The population-mean ICC/$\rho$ are near zero because a constant
predictor cannot rank vehicles \emph{within} a dimension; the large percentage
gains against it are therefore a change in kind (no ranking ability to substantial
agreement) more than a proportional improvement.
\end{minipage}
\end{table}

Performance varied by dimension (\Cref{fig:3seed_mae,fig:3seed_icc,fig:3seed_spearman}).
At $N{=}20$ the judge predicted ruggedness best (MAE $0.63$, ICC $0.78$) and overall
preference worst (MAE $0.87$--$0.94$, ICC $0.57$--$0.62$), with the three remaining
attributes in between. This ordering mirrors the human interrater reliability
reported in \Cref{sec:res-reliability}: the dimensions on which people agree most
are also the dimensions the judge predicts most reliably, and the holistic,
taste-driven preference judgment is hardest for both humans and model.

% =====================================================================
\subsection{The Distinct Behavior of Ruggedness}
\label{sec:disc-rugged-judge}
Ruggedness is anomalous across the judge experiments in a way that is consistent
with its anomalous behavior in the human data (\Cref{sec:res-reliability,fig:corr}).
Three observations stand out. First, as noted above, it is the single dimension the
zero-shot judge under-predicts: whereas the model reflexively inflates preference
and the other attributes, it \emph{under}-calls ruggedness, consistent with rugged
styling being a low-base-rate, minority aesthetic (\Cref{fig:dist}) that a
context-free model does not expect. Second, and most notably, ruggedness is the
best-agreed dimension at \emph{every} context size: its ICC is already $0.61$ at
zero-shot---when every other dimension sits between $0.18$ and $0.37$---and it rises
to $\approx 0.78$--$0.80$ by $N{=}20$, the highest of any dimension.

This produces an instructive dissociation between error and agreement. Absolute
error (MAE) and rank agreement (ICC/$\rho$) measure different things: MAE penalizes
calibration offsets, whereas ICC and $\rho$ reward getting the \emph{ordering} of a
person's vehicles right. Ruggedness scores well on agreement long before its error
is competitive because ruggedness is visually the most separable attribute---high
ground clearance, body cladding, and an upright stance make a vehicle
unambiguously rugged or not---so the model orders vehicles along this axis
easily even at zero-shot, and only needs context to correct its \emph{absolute}
level (the negative bias). Preference is the mirror image: it is the most holistic
and idiosyncratic judgment, so it carries both the largest zero-shot bias
($+0.85$) and the weakest agreement, and it benefits the most, proportionally,
from seeing how a particular person actually scores vehicles. That the judge ends
up predicting the crowd's most-agreed dimension (ruggedness) most reliably, and its
most-contested dimension (preference) least reliably, indicates that the judge is
recovering the same structure present in the human ratings rather than an artifact
of the model.
